\documentclass[ekobook.tex]{subfiles}

\begin{document}\IfDarkModeTF{\definecolor{WPP}{HTML}{DAB852}}{\definecolor{WPP}{HTML}{663800}}
\pageornament{31}{WPP}{3cm}
\vspace*{\fill}
\begin{center}
\begin{minipage}[c]{0.70\textwidth}
{\HUGE\color{WPP}\textbf{Maturitní otázky \colorbox{WPP}{\color{DP}Ekonomie}}}\\
\vspace{0em}%
{\color{WPP}\pgfornament[width=\textwidth]{88}}%
\vspace{0em}\par
\begin{center}
{\large Seznam dvaceti otázek a stručných odpovědí.}\par
{\large Vypracoval \href{https://github.com/ScamanderWayne}{\color{WPP}Libor Halík}}.
\end{center}
\end{minipage}
\end{center}
\vspace*{\fill}
\vspace*{\fill}
\thispagestyle{empty}
\newpage
\vspace*{\fill}    % odtlačí vše dolů

\hfill\begin{minipage}{0.6\textwidth}
\raggedleft               % zarovná text uvnitř minipage doprava
\parindent=0pt            % žádné odsazení odstavců (v tiráži to vypadá lépe)
Zdravím, tady \textbl{Libor}, někdy také \textbl{Scamander}.
Jsem autorem tohoto zpracování přípravy na maturitní zkoušku.
Zpracovával jsem to v rámci \texthl{školního roku 2025/26}, používal jsem školní materiály a aktuální znění dle zdrojů z internetu.
Tedy procentní částky, lhůty a další informace, by měli být přesné, dokud je vláda znovu nezmění.

\bigskip
Pokud si všimnete jakékoli hrubky.
Ať už gramatické, či faktické (například změna znění zákonů, kvůli které se \wemph{znovu} změní DPH), neváhejte se ozvat na mail \texthl{libor.halik.sw@gmail.com},
či můžete vytvořit ticket na \href{https://github.com/ScamanderWayne/WayneTeX/issues/new/choose}{\color{WPP}WayneTeX GitHub repository}.

\bigskip
Jestli jste technický nadšenec jako já, můžete celý projekt forknout.
Doporučuji kompilovat přes Lua\LaTeX .
\end{minipage}
\wpage
%\hfill\begin{minipage}[t]{0.15\textwidth}\raggedleft\wloss\end{minipage}%
\renewcommand\contentsname{{\color{WPP}Maturitní otázky{ }}\colorbox{WPP}{\color{DP}Ekonomie}}
\hypertarget{contents}{}
\tableofcontents
\thispagestyle{empty}
\begin{tikzpicture}[remember picture, overlay]
    % First node (Node A)
    \node[anchor=north east] (nodeA) at (current page.north east) {};
    % Second node (Node B) positioned -2em x and -2em y from Node A
    \node[below left=2.5em and 2em of nodeA] (nodeB) {\wloss};
\end{tikzpicture}
\wpage
\end{document}